% Explanation certificate for heart-s5-c2 (genome 57c311c6-56f5-4c5d-a475-46433ef58bf2).
% Generated by explaneat export (schema 1.4.0) from the JSON record alone;
% content sha256 e5293773fce87faec815a166111f195edf55b68bde14ba046b51556ab442dcea.
% Drafter recorded in the JSON: ai_agent (claude). Do not edit by hand: re-render with
%   uv run python -m explaneat export render --from heart-s5-c2.json
% TikZ modules are input from \explaneatcertdir{<label>} (empty by default: same
% directory; index.tex sets it to the genome's subdirectory).
\providecommand{\explaneatcertdir}[1]{}
\begin{drafted}[appendix: heart-s5-c2 explanation certificate]
\section{Explanation certificate: heart-s5-c2}\label{app:heart-s5-c2:sec}
This section is the explanation certificate of the heart-s5-c2 network. It is rendered from the canonical JSON record \texttt{heart-s5-c2.json} (schema 1.4.0, SHA-256 \texttt{e5293773fce87fae}\ldots), which also holds every evidence payload, the full graphs with weights, and the explanation record. \Cref{app:heart-s5-c2:summary} summarises the network and its coverage.

\begin{table}[htbp]
\centering\small
\begin{tabularx}{\linewidth}{@{}lX@{}}
\toprule
Dataset & Heart Disease (binary, prepared) \\
Experiment & 10seed-heart-s5, seed 5 \\
Genome & \texttt{57c311c6-56f5-4c5d-a475-46433ef58bf2} (NEAT key 9858, generation 39) \\
Explanation & \texttt{dabdb4bb-e227-46eb-a995-d4123373f234} (39 operations) \\
Train AUC / accuracy & 0.9389 / 0.8802 (n = 242) \\
Test AUC / accuracy & 0.8922 / 0.8197 (n = 61) \\
Output activation & sigmoid \\
Phenotype & 32 nodes (26 inputs, 5 hidden, 1 output), 35 enabled edges \\
Explained model state & 41 nodes, 37 enabled edges \\
Refactoring operations & add\_\allowbreak{}identity\_\allowbreak{}node 2, split\_\allowbreak{}node 5 \\
Function preservation & max $|\Delta\,\mathrm{output}|$ = $2.22\times 10^{-16}$ on the training split \\
Annotations & 6 (5 leaf, 1 composition) \\
Structural coverage & 1 (40 of 40 candidate nodes; 5 Null) \\
Compositional coverage & 1 \\
Uncovered nodes / edges & 0 / 0 \\
Missing compositions & 0 \\
Full coverage & yes \\
Drafted by & AI-assisted draft \\
Verified by & unverified \\
Record hash (SHA-256) & \texttt{e5293773fce87fae}\ldots \\
\bottomrule
\end{tabularx}
\caption{Explanation certificate summary: heart-s5-c2.}
\label{app:heart-s5-c2:summary}
\end{table}

\paragraph{Drafter's summary.} Mostly logistic regression: 17 direct edges (logit SD 3.19; train AUC 0.933 alone) plus a small ischaemia block (logit SD 0.57 for all hidden readouts) to 0.939. An ischaemia gate ReLU(0.39 + 2.54 oldpeak + 0.99 restecg:0 - 0.88 thal:4) opens above an ST-depression threshold that is lower with a normal resting ECG and raises risk directly and through a gate that passes the score on minus a resting-blood-pressure offset (more risk when BP is low); a near-linear node takes back part of the ischaemia score and adds a fixed-defect term. Two small indicators add risk for LVH without exercise angina and (negligibly) for high fasting sugar. Same ischaemia-gate -$>$ resting-BP-gate core as the 10-seed heart-s5 record (key 9330) and the sibling lead heart-s5-c1, but without their age gate.

\paragraph{What resisted explanation.} Nothing resisted coverage. The ischaemia gate 204 feeds two downstream gates (290, 1416) as well as the output; the two links were intercepted with identity nodes (identity\_\allowbreak{}290, identity\_\allowbreak{}1416) so that each hidden node is its own leaf, and the root composes all five. Counter-directional signs reported, not explained: restecg:0 (normal ECG) raises the ischaemia score (its direct edge is -0.16), trestbps enters the risk-raising BP gate negatively (low resting BP opens it), thal:6 (fixed defect) has a direct -0.41 edge (partly offset by +0.07/SD through 1416), fbs:0 +0.26 directly so high fasting sugar lowers the logit overall despite the +0.07 indicator. thal:4 shuts the ischaemia gate but has only 2 training patients. The damping node 1416 and the fasting-sugar node 3884 are effectively linear (affine R2 0.99 and 1.0); the LVH indicator and fasting-sugar indicator barely affect ranking (readout ablation $<$= -0.0003 AUC).

\subsection{Plain-language summary}\label{app:heart-s5-c2:plain}
\noindent{\small\textit{AI-assisted draft for review (unverified). Each sentence cites the claims (C01, C02, \ldots) that back it; every cited claim was computed on the model and the training split and passed its automatic checks. Magnitudes are changes in predicted risk for a typical patient; confidence words follow the test-split check and n. The text describes what the model does, not the world.}}\par\smallskip
\paragraph{Summary (drafted).} For a typical patient (a 56-year-old man with no chest pain symptoms and median test results) the model predicts a risk of heart disease of about 20\%; it adds up most of its inputs independently, and five parts add to that sum: a signal from ST depression on exercise and the resting ECG that acts on 109 of 242 training patients, a blood pressure step and a damping term that both read that signal, an extra term for left ventricular hypertrophy without chest pain on exercise, and a very small term for high fasting blood sugar.\,{\scriptsize\textcolor{NavyBlue}{[C34, C42, C53, C61]}} These add up independently, and these lower the risk: being a woman takes a typical patient's predicted risk from about 20\% to about 3\%; through this sum alone typical angina lowers it to about 2\%, non-anginal chest pain to about 4\%, fasting blood sugar above 120 mg/dL to about 10\% and atypical angina to about 11\%; and a higher maximum heart rate lowers it steadily, from about 26\% at 118 to about 16\% at 175 beats per minute.\,{\scriptsize\textcolor{NavyBlue}{[C08, C06, C12, C18, C20, C19]}} Each extra vessel coloured by fluoroscopy adds about the same, from about 20\% at 0 to about 87\% at 2 vessels, flattening out near the top because the predicted risk cannot go much higher; a flat rather than upsloping ST segment raises it to about 61\%, higher cholesterol raises it steadily from about 13\% at 193 mg/dL to about 34\% at 308 mg/dL, and older age raises it steadily from about 17\% at 42 to about 22\% at 66 years.\,{\scriptsize\textcolor{NavyBlue}{[C01, C02, C04, C27]}} Through this sum alone, higher resting blood pressure raises it steadily, from about 12\% at 110 mm Hg to about 33\% at 155 mm Hg, a thallium result other than a normal scan or a fixed defect raises it to about 36\%, chest pain brought on by exercise raises it to about 33\%, and a resting ECG other than normal raises it to about 25\%; all four also act through the parts below.\,{\scriptsize\textcolor{NavyBlue}{[C05, C09, C14, C26]}} Some rarer results also count (the model has little data on this, n = 15 to 17): through this sum alone a fixed defect on the thallium scan lowers the risk to about 9\%, and a downsloping ST segment raises it to about 27\%.\,{\scriptsize\textcolor{NavyBlue}{[C17, C24]}} This part reads ST depression on exercise, the resting ECG and the thallium scan, and acts on 109 of 242 training patients, 65 of whom had heart disease; for the typical patient it acts on, the predicted risk is about 76\% without it and about 89\% with it.\,{\scriptsize\textcolor{NavyBlue}{[C42]}} Through it, more ST depression adds to the predicted risk only once it rises above about 1.3 mm (between 0.4 and 1.3 mm for different patients), rising in an uneven way: for a typical patient from about 16\% at 0 mm to about 50\% at 2.8 mm.\,{\scriptsize\textcolor{NavyBlue}{[C35, C41]}} ST depression counts more with a normal resting ECG (from about 16\% to about 50\%) than with any other resting ECG (from about 25\% to about 48\%), in an uneven way.\,{\scriptsize\textcolor{NavyBlue}{[C39]}} Through it, a resting ECG other than normal lowers a typical patient's risk from about 20\% to about 16\%, and a thallium result coded 4 shuts the signal (the model has little data on this, n = 2).\,{\scriptsize\textcolor{NavyBlue}{[C36, C37]}} This part reads that signal together with resting blood pressure, chest pain on exercise and chest pain type, and acts on 145 of 242 training patients, 73 of whom had heart disease; for the typical patient it acts on, the predicted risk is about 51\% without it and about 68\% with it.\,{\scriptsize\textcolor{NavyBlue}{[C53]}} Through it, higher resting blood pressure lowers the risk up to about 125 mm Hg (between 122 and 165 mm Hg for different patients), then no further, in an uneven way: for a typical patient from about 23\% at 110 mm Hg to about 17\% at 155 mm Hg.\,{\scriptsize\textcolor{NavyBlue}{[C44, C52]}} Through it, chest pain on exercise lowers a typical patient's risk from about 20\% to about 18\%, and non-anginal chest pain to about 19\%.\,{\scriptsize\textcolor{NavyBlue}{[C45, C46]}} This part reads the same signal and the thallium scan, and acts on 234 of 242 training patients, 102 of whom had heart disease; for the typical patient it acts on, the predicted risk is about 24\% without it and about 20\% with it.\,{\scriptsize\textcolor{NavyBlue}{[C56]}} Through it, a fixed defect on the thallium scan raises a typical patient's risk from about 20\% to about 23\% (the model has little data on this, n = 15).\,{\scriptsize\textcolor{NavyBlue}{[C55]}} This part acts on 74 of 242 training patients, nearly all with left ventricular hypertrophy on the resting ECG and no chest pain on exercise, 30 of whom had heart disease; for the typical patient it acts on, the predicted risk is about 25\% without it and about 30\% with it.\,{\scriptsize\textcolor{NavyBlue}{[C61]}} Through it, left ventricular hypertrophy raises a typical patient's risk from about 20\% to about 24\%, and for the typical patient it acts on, chest pain on exercise lowers it from about 30\% to about 25\%.\,{\scriptsize\textcolor{NavyBlue}{[C59, C58]}} Across the model, chest pain on exercise raises a typical patient's risk from about 20\% to about 31\% without left ventricular hypertrophy, and from about 25\% to about 35\% with it.\,{\scriptsize\textcolor{NavyBlue}{[C60]}} This part acts on 33 of 242 training patients, nearly all with fasting blood sugar above 120 mg/dL, 15 of whom had heart disease; for the typical patient it acts on, the predicted risk is about 24\% without it and about 25\% with it.\,{\scriptsize\textcolor{NavyBlue}{[C65]}} It barely changes the predictions, by at most 2 percentage points.\,{\scriptsize\textcolor{NavyBlue}{[C66]}}
\paragraph{The whole model.} The model adds up its inputs into one score: the vessels coloured (each extra vessel adding about the same), the ST segment, sex and chest pain type each move the predicted risk on their own, and five parts add to that sum: a signal from ST depression and the resting ECG, a blood pressure step and a damping term that read it, a hypertrophy term and a fasting sugar term.\,{\scriptsize\textcolor{NavyBlue}{[C34, C01, C02, C08, C06]}} \textit{Evidence.} It is made up of the inputs that feed the score directly and five parts (each explained below); the first part feeds the second and the third, and the score is their sum, which was checked on all 242 training patients.\,{\scriptsize\textcolor{NavyBlue}{[C34]}} ST depression on exercise does not feed the score directly, so it acts only through these parts.\,{\scriptsize\textcolor{NavyBlue}{[C34, C35]}} Changing one input at a time across the whole model, more ST depression raises a typical patient's predicted risk in an uneven way, from about 16\% at 0 mm to about 50\% at 2.8 mm, and higher resting blood pressure raises it in an uneven way, from about 15\% at 110 mm Hg to about 28\% at 155 mm Hg; through the sum alone it rises steadily, from about 12\% to about 33\%.\,{\scriptsize\textcolor{NavyBlue}{[C03, C13, C05]}} Across the whole model, chest pain on exercise raises it to about 31\%, non-anginal chest pain lowers it to about 4\% and fasting blood sugar above 120 mg/dL to about 11\%.\,{\scriptsize\textcolor{NavyBlue}{[C15, C11, C22]}} Across the whole model, a fixed defect on the thallium scan lowers it to about 11\% (the model has little data on this, n = 15).\,{\scriptsize\textcolor{NavyBlue}{[C23]}}
\paragraph{ST depression signal.} This part raises the predicted risk with ST depression on exercise once it rises above about 1.3 mm, in an uneven way, and more so with a normal resting ECG.\,{\scriptsize\textcolor{NavyBlue}{[C35, C41, C39]}} \textit{Evidence.} It acts on 109 of 242 training patients, 65 of whom had heart disease; for the typical patient it acts on, the predicted risk is about 76\% without it and about 89\% with it.\,{\scriptsize\textcolor{NavyBlue}{[C42]}} Through it, more ST depression adds to the risk only once it rises above about 1.3 mm (between 0.4 and 1.3 mm for different patients), rising in an uneven way: for a typical patient from about 16\% at 0 mm to about 50\% at 2.8 mm.\,{\scriptsize\textcolor{NavyBlue}{[C35, C41]}} ST depression counts more with a normal resting ECG (from about 16\% to about 50\%) than with any other resting ECG (from about 25\% to about 48\%), in an uneven way.\,{\scriptsize\textcolor{NavyBlue}{[C39]}} Through it, a resting ECG other than normal lowers a typical patient's risk from about 20\% to about 16\%, and a thallium result coded 4 shuts the signal (the model has little data on this, n = 2).\,{\scriptsize\textcolor{NavyBlue}{[C36, C37]}} Its formula reproduces the part exactly on all 242 training patients.\,{\scriptsize\textcolor{NavyBlue}{[C43]}}
\paragraph{Blood pressure step.} This part passes the ST depression signal on to the predicted risk, more for patients with lower resting blood pressure: through it, higher resting blood pressure lowers the risk in an uneven way up to about 125 mm Hg (between 122 mm Hg and 165 mm Hg for different patients), then no further.\,{\scriptsize\textcolor{NavyBlue}{[C53, C44, C52]}} \textit{Evidence.} It acts on 145 of 242 training patients, 73 of whom had heart disease; for the typical patient it acts on, the predicted risk is about 51\% without it and about 68\% with it.\,{\scriptsize\textcolor{NavyBlue}{[C53]}} Through it, higher resting blood pressure lowers the risk up to about 125 mm Hg (between 122 and 165 mm Hg for different patients), then no further, in an uneven way: for a typical patient from about 23\% at 110 mm Hg to about 17\% at 155 mm Hg.\,{\scriptsize\textcolor{NavyBlue}{[C44, C52]}} Through it, chest pain on exercise lowers a typical patient's risk from about 20\% to about 18\%, and non-anginal chest pain to about 19\%.\,{\scriptsize\textcolor{NavyBlue}{[C45, C46]}} Its formula reproduces the part exactly on all 242 training patients.\,{\scriptsize\textcolor{NavyBlue}{[C54]}}
\paragraph{Damping term.} This part takes back a little of the ST depression signal for nearly every patient and adds a little for a fixed defect on the thallium scan (the model has little data on this, n = 15).\,{\scriptsize\textcolor{NavyBlue}{[C56, C55]}} \textit{Evidence.} It acts on 234 of 242 training patients, 102 of whom had heart disease; for the typical patient it acts on, the predicted risk is about 24\% without it and about 20\% with it.\,{\scriptsize\textcolor{NavyBlue}{[C56]}} Through it, a fixed defect on the thallium scan raises a typical patient's risk from about 20\% to about 23\% (the model has little data on this, n = 15).\,{\scriptsize\textcolor{NavyBlue}{[C55]}} Its formula reproduces the part exactly on all 242 training patients.\,{\scriptsize\textcolor{NavyBlue}{[C57]}}
\paragraph{Hypertrophy without exercise pain.} This part adds a little to the predicted risk for patients with left ventricular hypertrophy on the resting ECG and no chest pain on exercise.\,{\scriptsize\textcolor{NavyBlue}{[C61, C59]}} \textit{Evidence.} It acts on 74 of 242 training patients, nearly all with left ventricular hypertrophy and no chest pain on exercise, 30 of whom had heart disease; for the typical patient it acts on, the predicted risk is about 25\% without it and about 30\% with it.\,{\scriptsize\textcolor{NavyBlue}{[C61]}} Through it, left ventricular hypertrophy raises a typical patient's risk from about 20\% to about 24\%, and for the typical patient it acts on, chest pain on exercise lowers it from about 30\% to about 25\%.\,{\scriptsize\textcolor{NavyBlue}{[C59, C58]}} Across the model, chest pain on exercise raises a typical patient's risk from about 20\% to about 31\% without left ventricular hypertrophy, and from about 25\% to about 35\% with it.\,{\scriptsize\textcolor{NavyBlue}{[C60]}} Its formula reproduces the part exactly on all 242 training patients.\,{\scriptsize\textcolor{NavyBlue}{[C63]}}
\paragraph{High fasting sugar term.} This part adds a very small amount to the predicted risk for patients with fasting blood sugar above 120 mg/dL.\,{\scriptsize\textcolor{NavyBlue}{[C65, C64]}} \textit{Evidence.} It acts on 33 of 242 training patients, nearly all with fasting blood sugar above 120 mg/dL, 15 of whom had heart disease; for the typical patient it acts on, the predicted risk is about 24\% without it and about 25\% with it.\,{\scriptsize\textcolor{NavyBlue}{[C65]}} Through it, fasting blood sugar above 120 mg/dL raises a typical patient's risk from about 20\% to about 21\%, and it changes the predictions by at most 2 percentage points.\,{\scriptsize\textcolor{NavyBlue}{[C64, C66]}} Its formula reproduces the part exactly on all 242 training patients.\,{\scriptsize\textcolor{NavyBlue}{[C67]}}
\paragraph{Worth checking against what is known.} Worth checking: the model treats fasting blood sugar above 120 mg/dL as lowering the risk of heart disease (for a typical patient from about 20\% to about 11\%).\,{\scriptsize\textcolor{NavyBlue}{[C22]}} Worth checking: patients with typical angina get the lowest predicted risk of all chest pain types, below patients with no chest pain symptoms.\,{\scriptsize\textcolor{NavyBlue}{[C07, C11, C21]}} Worth checking: the model treats a fixed defect on the thallium scan as lowering the risk compared with a normal scan, from only 15 training patients.\,{\scriptsize\textcolor{NavyBlue}{[C23]}} Worth checking: ST depression on exercise raises the predicted risk more when the resting ECG is normal than when it is not.\,{\scriptsize\textcolor{NavyBlue}{[C39]}} Worth checking: through the blood pressure step, higher resting blood pressure lowers the predicted risk, against its rise through the sum.\,{\scriptsize\textcolor{NavyBlue}{[C44, C05]}} Worth checking: for the 2 training patients whose thallium result is coded 4, a code the dataset documentation does not define, the ST depression signal is shut.\,{\scriptsize\textcolor{NavyBlue}{[C37]}}
\paragraph{Claims.}
\par{\tiny\noindent
\begin{tabularx}{\linewidth}{@{}l>{\raggedright\arraybackslash}X>{\raggedright\arraybackslash}p{0.11\linewidth}>{\raggedright\arraybackslash}p{0.16\linewidth}>{\raggedright\arraybackslash}p{0.05\linewidth}>{\raggedright\arraybackslash}p{0.15\linewidth}l>{\raggedright\arraybackslash}p{0.08\linewidth}@{}}
\toprule
Claim & What is checked & Risk, typical patient & Estimate (log-odds; OR); range across patients & n & Test check & Verified & Words \\
\midrule
C01 & direction: number of vessels coloured by fluoroscopy: 0 vessels to 2 vessels (whole model) \newline\textit{shape: steady (equal steps); 100\% of patients same; verified; risk saturates near 100\%} & 20\% to 87\% & log-odds 3.30, OR 27.2; range across patients 3.30 to 3.30 & 242 & train; test unclear (AUC gain 0.025 [-0.049, 0.102], n=61) & \textcolor{OliveGreen}{yes} & on the training data only \\
C02 & direction: slope of the ST segment at peak exercise: an upsloping ST segment to a flat ST segment (whole model) \newline\textit{shape: categorical; highest to lowest: a flat ST segment, a downsloping ST segment, an upsloping ST segment} & 20\% to 61\% & log-odds 1.84, OR 6.32; range across patients 1.84 to 1.84 & 115, 110 & train; test disagrees (AUC gain -0.017 [-0.061, 0.016], n=61) & \textcolor{OliveGreen}{yes} & on the training data only \\
C03 & direction: ST depression on exercise: 0.0 mm to 2.8 mm (whole model) \newline\textit{shape: hinge at 0.4 mm; 70\% of patients same; fallback} & 16\% to 50\% & log-odds 1.22, OR 3.40; range across patients 0.77 to 1.69 & 242 & train; test unclear (AUC gain 0.009 [-0.008, 0.033], n=61) & \textcolor{OliveGreen}{yes} & on the training data only \\
C04 & direction: serum cholesterol: 193 mg/dL (5.0 mmol/L) to 308 mg/dL (8.0 mmol/L) (whole model) \newline\textit{shape: steady; 100\% of patients same; verified} & 13\% to 34\% & log-odds 1.25, OR 3.48; range across patients 1.25 to 1.25 & 242 & train; test disagrees (AUC gain -0.028 [-0.071, 0.002], n=61) & \textcolor{OliveGreen}{yes} & on the training data only \\
C05 & direction: resting blood pressure: 110 mm Hg to 155 mm Hg (direct contribution only, not through any part) \newline\textit{shape: steady; 100\% of patients same; verified} & 12\% to 33\% & log-odds 1.23, OR 3.43; range across patients 1.23 to 1.23 & 242 & train; test unclear (AUC gain 0.006 [-0.024, 0.037], n=61) & \textcolor{OliveGreen}{yes} & on the training data only \\
C06 & direction: chest pain type: no chest pain (asymptomatic) to typical angina (direct contribution only, not through any part) \newline\textit{shape: categorical; highest to lowest: no chest pain (asymptomatic), atypical angina, non-anginal chest pain, typical angina} & 20\% to 2\% & log-odds -2.42, OR 0.09; range across patients -2.42 to -2.42 & 108, 22 & train; test agrees (AUC gain 0.032 [0.008, 0.062], n=61) & \textcolor{OliveGreen}{yes} & clearly \\
C07 & direction: chest pain type: no chest pain (asymptomatic) to typical angina (whole model) \newline\textit{shape: categorical; highest to lowest: no chest pain (asymptomatic), atypical angina, non-anginal chest pain, typical angina} & 20\% to 2\% & log-odds -2.42, OR 0.09; range across patients -2.42 to -2.42 & 108, 22 & train; test agrees (AUC gain 0.032 [0.008, 0.062], n=61) & \textcolor{OliveGreen}{yes} & clearly \\
C08 & direction: sex: male to female (whole model) \newline\textit{shape: categorical; highest to lowest: male, female} & 20\% to 3\% & log-odds -2.24, OR 0.11; range across patients -2.24 to -2.24 & 159, 83 & train; test unclear (AUC gain 0.022 [-0.011, 0.065], n=61) & \textcolor{OliveGreen}{yes} & on the training data only \\
C09 & direction: thallium stress test result: a normal thallium scan to thallium stress test result other than a normal thallium scan, a fixed defect on the thallium scan (direct contribution only, not through any part) \newline\textit{shape: categorical; highest to lowest: thallium stress test result other than a normal thallium scan, a fixed defect on the thallium scan, a normal thallium scan, a fixed defect on the thallium scan} & 20\% to 36\% & log-odds 0.85, OR 2.34; range across patients 0.85 to 0.85 & 135, 92 & train; test unclear (AUC gain 0.013 [-0.011, 0.041], n=61) & \textcolor{OliveGreen}{yes} & on the training data only \\
C10 & direction: thallium stress test result: a normal thallium scan to a reversible defect on the thallium scan (whole model) \newline\textit{shape: categorical; highest to lowest: a reversible defect on the thallium scan, thallium stress test result coded 4 (meaning not recorded), a normal thallium scan, a fixed defect on the thallium scan} & 20\% to 36\% & log-odds 0.85, OR 2.34; range across patients 0.85 to 0.85 & 135, 90 & train; test unclear (AUC gain 0.014 [-0.009, 0.042], n=61) & \textcolor{OliveGreen}{yes} & on the training data only \\
C11 & direction: chest pain type: no chest pain (asymptomatic) to non-anginal chest pain (whole model) \newline\textit{shape: categorical; highest to lowest: no chest pain (asymptomatic), atypical angina, non-anginal chest pain, typical angina} & 20\% to 4\% & log-odds -1.72, OR 0.18; range across patients -1.74 to -1.70 & 108, 70 & train; test agrees (AUC gain 0.032 [0.008, 0.062], n=61) & \textcolor{OliveGreen}{yes} & clearly \\
C12 & direction: chest pain type: no chest pain (asymptomatic) to non-anginal chest pain (direct contribution only, not through any part) \newline\textit{shape: categorical; highest to lowest: no chest pain (asymptomatic), atypical angina, non-anginal chest pain, typical angina} & 20\% to 4\% & log-odds -1.70, OR 0.18; range across patients -1.70 to -1.70 & 108, 70 & train; test agrees (AUC gain 0.032 [0.008, 0.062], n=61) & \textcolor{OliveGreen}{yes} & clearly \\
\bottomrule
\end{tabularx}}
\par{\tiny\noindent
\begin{tabularx}{\linewidth}{@{}l>{\raggedright\arraybackslash}X>{\raggedright\arraybackslash}p{0.11\linewidth}>{\raggedright\arraybackslash}p{0.16\linewidth}>{\raggedright\arraybackslash}p{0.05\linewidth}>{\raggedright\arraybackslash}p{0.15\linewidth}l>{\raggedright\arraybackslash}p{0.08\linewidth}@{}}
\toprule
Claim & What is checked & Risk, typical patient & Estimate (log-odds; OR); range across patients & n & Test check & Verified & Words \\
\midrule
C13 & direction: resting blood pressure: 110 mm Hg to 155 mm Hg (whole model) \newline\textit{shape: accelerating; 75\% of patients same; fallback} & 15\% to 28\% & log-odds 0.97, OR 2.63; range across patients 0.75 to 1.15 & 242 & train; test unclear (AUC gain 0.005 [-0.017, 0.031], n=61) & \textcolor{OliveGreen}{yes} & on the training data only \\
C14 & direction: exercise-induced angina: no chest pain on exercise to chest pain brought on by exercise (direct contribution only, not through any part) \newline\textit{shape: categorical; highest to lowest: chest pain brought on by exercise, no chest pain on exercise} & 20\% to 33\% & log-odds 0.71, OR 2.03; range across patients 0.71 to 0.71 & 160, 82 & train; test agrees (AUC gain 0.012 [0.001, 0.030], n=61) & \textcolor{OliveGreen}{yes} & clearly \\
C15 & direction: exercise-induced angina: no chest pain on exercise to chest pain brought on by exercise (whole model) \newline\textit{shape: categorical; highest to lowest: chest pain brought on by exercise, no chest pain on exercise} & 20\% to 31\% & log-odds 0.55, OR 1.73; range across patients 0.39 to 0.71 & 160, 82 & train; test unclear (AUC gain 0.008 [0.000, 0.022], n=61) & \textcolor{OliveGreen}{yes} & on the training data only \\
C16 & direction: thallium stress test result: a normal thallium scan to thallium stress test result coded 4 (meaning not recorded) (whole model) \newline\textit{shape: categorical; highest to lowest: a reversible defect on the thallium scan, thallium stress test result coded 4 (meaning not recorded), a normal thallium scan, a fixed defect on the thallium scan} & 20\% to 31\% & log-odds 0.52, OR 1.69; range across patients -0.22 to 0.85 & 135, 2 (small n) & train; test unclear (AUC gain 0.014 [-0.009, 0.042], n=61) & \textcolor{BrickRed}{no} & the model has little data on this; on the training data only \\
C17 & direction: thallium stress test result: a normal thallium scan to a fixed defect on the thallium scan (direct contribution only, not through any part) \newline\textit{shape: categorical; highest to lowest: thallium stress test result other than a normal thallium scan, a fixed defect on the thallium scan, a normal thallium scan, a fixed defect on the thallium scan} & 20\% to 9\% & log-odds -0.86, OR 0.42; range across patients -0.86 to -0.86 & 135, 15 (small n) & train; test unclear (AUC gain 0.013 [-0.011, 0.041], n=61) & \textcolor{OliveGreen}{yes} & the model has little data on this; on the training data only \\
C18 & direction: fasting blood sugar: fasting blood sugar of 120 mg/dL or less to fasting blood sugar above 120 mg/dL (direct contribution only, not through any part) \newline\textit{shape: categorical; highest to lowest: fasting blood sugar of 120 mg/dL or less, fasting blood sugar above 120 mg/dL} & 20\% to 10\% & log-odds -0.77, OR 0.46; range across patients -0.77 to -0.77 & 209, 33 & train; test disagrees (AUC gain -0.004 [-0.020, 0.009], n=61) & \textcolor{OliveGreen}{yes} & on the training data only \\
C19 & direction: maximum heart rate: 118 beats per minute to 175 beats per minute (whole model) \newline\textit{shape: steady; 100\% of patients same; verified} & 26\% to 16\% & log-odds -0.57, OR 0.56; range across patients -0.57 to -0.57 & 242 & train; test unclear (AUC gain 0.006 [0.000, 0.022], n=61) & \textcolor{OliveGreen}{yes} & on the training data only \\
C20 & direction: chest pain type: no chest pain (asymptomatic) to atypical angina (direct contribution only, not through any part) \newline\textit{shape: categorical; highest to lowest: no chest pain (asymptomatic), atypical angina, non-anginal chest pain, typical angina} & 20\% to 11\% & log-odds -0.72, OR 0.49; range across patients -0.72 to -0.72 & 108, 42 & train; test agrees (AUC gain 0.032 [0.008, 0.062], n=61) & \textcolor{OliveGreen}{yes} & clearly \\
C21 & direction: chest pain type: no chest pain (asymptomatic) to atypical angina (whole model) \newline\textit{shape: categorical; highest to lowest: no chest pain (asymptomatic), atypical angina, non-anginal chest pain, typical angina} & 20\% to 11\% & log-odds -0.72, OR 0.49; range across patients -0.72 to -0.72 & 108, 42 & train; test agrees (AUC gain 0.032 [0.008, 0.062], n=61) & \textcolor{OliveGreen}{yes} & clearly \\
C22 & direction: fasting blood sugar: fasting blood sugar of 120 mg/dL or less to fasting blood sugar above 120 mg/dL (whole model) \newline\textit{shape: categorical; highest to lowest: fasting blood sugar of 120 mg/dL or less, fasting blood sugar above 120 mg/dL} & 20\% to 11\% & log-odds -0.69, OR 0.50; range across patients -0.69 to -0.69 & 209, 33 & train; test disagrees (AUC gain -0.003 [-0.019, 0.009], n=61) & \textcolor{OliveGreen}{yes} & on the training data only \\
C23 & direction: thallium stress test result: a normal thallium scan to a fixed defect on the thallium scan (whole model) \newline\textit{shape: categorical; highest to lowest: a reversible defect on the thallium scan, thallium stress test result coded 4 (meaning not recorded), a normal thallium scan, a fixed defect on the thallium scan} & 20\% to 11\% & log-odds -0.69, OR 0.50; range across patients -0.75 to -0.58 & 135, 15 (small n) & train; test unclear (AUC gain 0.014 [-0.009, 0.042], n=61) & \textcolor{OliveGreen}{yes} & the model has little data on this; on the training data only \\
C24 & direction: slope of the ST segment at peak exercise: an upsloping ST segment to a downsloping ST segment (whole model) \newline\textit{shape: categorical; highest to lowest: a flat ST segment, a downsloping ST segment, an upsloping ST segment} & 20\% to 27\% & log-odds 0.40, OR 1.48; range across patients 0.40 to 0.40 & 115, 17 (small n) & train; test disagrees (AUC gain -0.017 [-0.061, 0.016], n=61) & \textcolor{OliveGreen}{yes} & the model has little data on this; on the training data only \\
\bottomrule
\end{tabularx}}
\par{\tiny\noindent
\begin{tabularx}{\linewidth}{@{}l>{\raggedright\arraybackslash}X>{\raggedright\arraybackslash}p{0.11\linewidth}>{\raggedright\arraybackslash}p{0.16\linewidth}>{\raggedright\arraybackslash}p{0.05\linewidth}>{\raggedright\arraybackslash}p{0.15\linewidth}l>{\raggedright\arraybackslash}p{0.08\linewidth}@{}}
\toprule
Claim & What is checked & Risk, typical patient & Estimate (log-odds; OR); range across patients & n & Test check & Verified & Words \\
\midrule
C25 & direction: resting ECG: a normal resting ECG to signs of left ventricular hypertrophy on the resting ECG (whole model) \newline\textit{shape: categorical; highest to lowest: signs of left ventricular hypertrophy on the resting ECG, an ST-T wave abnormality on the resting ECG, a normal resting ECG} & 20\% to 25\% & log-odds 0.21, OR 1.24; range across patients -0.32 to 0.55 & 121, 119 & train; test disagrees (AUC gain -0.003 [-0.023, 0.013], n=61) & \textcolor{BrickRed}{no} & on the training data only \\
C26 & direction: resting ECG: a normal resting ECG to resting ECG other than a normal resting ECG (direct contribution only, not through any part) \newline\textit{shape: categorical; highest to lowest: resting ECG other than a normal resting ECG, a normal resting ECG} & 20\% to 25\% & log-odds 0.32, OR 1.37; range across patients 0.32 to 0.32 & 121, 121 & train; test unclear (AUC gain 0.000 [-0.016, 0.015], n=61) & \textcolor{OliveGreen}{yes} & on the training data only \\
C27 & direction: age: 42 years to 66 years (whole model) \newline\textit{shape: steady; 100\% of patients same; verified} & 17\% to 22\% & log-odds 0.33, OR 1.39; range across patients 0.33 to 0.33 & 242 & train; test disagrees (AUC gain -0.005 [-0.023, 0.01], n=61) & \textcolor{OliveGreen}{yes} & on the training data only \\
C28 & direction: resting ECG: a normal resting ECG to an ST-T wave abnormality on the resting ECG (whole model) \newline\textit{shape: categorical; highest to lowest: signs of left ventricular hypertrophy on the resting ECG, an ST-T wave abnormality on the resting ECG, a normal resting ECG} & 20\% to 21\% & log-odds 0.06, OR 1.06; range across patients -0.32 to 0.32 & 121, 2 (small n) & train; test disagrees (AUC gain -0.003 [-0.023, 0.013], n=61) & \textcolor{BrickRed}{no} & slightly; the model has little data on this; on the training data only \\
C29 & magnitude: number of vessels coloured by fluoroscopy: 0 vessels to 2 vessels, typical patient \newline\textit{shape: steady (equal steps); 100\% of patients same; verified; risk saturates near 100\%} & 20\% to 87\% & log-odds 3.30, OR 27.2 & 242 & train; test unclear (AUC gain 0.025 [-0.049, 0.102], n=61) & \textcolor{OliveGreen}{yes} & on the training data only \\
C30 & magnitude: slope of the ST segment at peak exercise: an upsloping ST segment to a flat ST segment, typical patient \newline\textit{shape: categorical; highest to lowest: a flat ST segment, a downsloping ST segment, an upsloping ST segment} & 20\% to 61\% & log-odds 1.84, OR 6.32 & 115, 110 & train; test disagrees (AUC gain -0.017 [-0.061, 0.016], n=61) & \textcolor{OliveGreen}{yes} & on the training data only \\
C31 & magnitude: ST depression on exercise: 0.0 mm to 2.8 mm, typical patient \newline\textit{shape: hinge at 0.4 mm; 70\% of patients same; fallback} & 16\% to 50\% & log-odds 1.64, OR 5.15 & 242 & train; test unclear (AUC gain 0.009 [-0.008, 0.033], n=61) & \textcolor{OliveGreen}{yes} & on the training data only \\
C32 & magnitude: serum cholesterol: 193 mg/dL (5.0 mmol/L) to 308 mg/dL (8.0 mmol/L), typical patient \newline\textit{shape: steady; 100\% of patients same; verified} & 13\% to 34\% & log-odds 1.25, OR 3.48 & 242 & train; test disagrees (AUC gain -0.028 [-0.071, 0.002], n=61) & \textcolor{OliveGreen}{yes} & on the training data only \\
C33 & magnitude: chest pain type: no chest pain (asymptomatic) to typical angina, typical patient \newline\textit{shape: categorical; highest to lowest: no chest pain (asymptomatic), atypical angina, non-anginal chest pain, typical angina} & 20\% to 2\% & log-odds -2.42, OR 0.09 & 108, 22 & train; test agrees (AUC gain 0.032 [0.008, 0.062], n=61) & \textcolor{OliveGreen}{yes} & clearly \\
C34 & structure: made of ischaemia\_\allowbreak{}gate, bp\_\allowbreak{}offset\_\allowbreak{}gate, ischaemia\_\allowbreak{}damping\_\allowbreak{}thal6,\allowbreak{} lvh\_\allowbreak{}without\_\allowbreak{}exercise\_\allowbreak{}angina,\allowbreak{} high\_\allowbreak{}fasting\_\allowbreak{}sugar and inputs age, sex, cp, trestbps, chol, fbs, restecg, thalach, exang, slope, ca, thal; added; ann\_\allowbreak{}33 feeds ann\_\allowbreak{}34; ann\_\allowbreak{}33 feeds ann\_\allowbreak{}35 &  & log-odds n/a, OR n/a & 242 & train & \textcolor{OliveGreen}{yes} &  \\
C35 & direction: ST depression on exercise: 0.0 mm to 2.8 mm (through this part) \newline\textit{shape: hinge at 1.3 mm; 70\% of patients same; fallback; one point with C41 (median fold 1.3 mm, typical patient 0.4 mm, 10th-90th percentile 0.4 mm to 1.3 mm, wide)} & 16\% to 50\% & log-odds 1.22, OR 3.40; range across patients 0.77 to 1.69 & 242 & train; test unclear (AUC gain 0.009 [-0.008, 0.033], n=61) & \textcolor{OliveGreen}{yes} & on the training data only \\
C36 & direction: resting ECG: a normal resting ECG to resting ECG other than a normal resting ECG (through this part) \newline\textit{shape: categorical; highest to lowest: a normal resting ECG, resting ECG other than a normal resting ECG} & 20\% to 16\% & log-odds -0.26, OR 0.77; range across patients -0.64 to 0.00 & 121, 121 & train; test disagrees (AUC gain -0.002 [-0.016, 0.012], n=61) & \textcolor{OliveGreen}{yes} & slightly; on the training data only \\
\bottomrule
\end{tabularx}}
\par{\tiny\noindent
\begin{tabularx}{\linewidth}{@{}l>{\raggedright\arraybackslash}X>{\raggedright\arraybackslash}p{0.11\linewidth}>{\raggedright\arraybackslash}p{0.16\linewidth}>{\raggedright\arraybackslash}p{0.05\linewidth}>{\raggedright\arraybackslash}p{0.15\linewidth}l>{\raggedright\arraybackslash}p{0.08\linewidth}@{}}
\toprule
Claim & What is checked & Risk, typical patient & Estimate (log-odds; OR); range across patients & n & Test check & Verified & Words \\
\midrule
C37 & direction: thallium stress test result: thallium stress test result other than thallium stress test result coded 4 (meaning not recorded) to thallium stress test result coded 4 (meaning not recorded) (through this part) \newline\textit{shape: categorical; highest to lowest: thallium stress test result other than thallium stress test result coded 4 (meaning not recorded), thallium stress test result coded 4 (meaning not recorded)} & 20\% to 16\% & log-odds -0.33, OR 0.72; range across patients -1.07 to 0.00 & 240, 2 (small n) & train; test disagrees (AUC gain -0.001 [-0.006, 0.000], n=61) & \textcolor{OliveGreen}{yes} & slightly; the model has little data on this; on the training data only \\
C38 & interaction: effect of ST depression on exercise (0.0 mm to 2.8 mm) by thallium stress test result: thallium stress test result coded 4 (meaning not recorded) minus thallium stress test result other than thallium stress test result coded 4 (meaning not recorded); stronger in thallium stress test result other than thallium stress test result coded 4 (meaning not recorded) \newline\textit{shape: hinge at 0.4 mm; 74\% of patients same; fallback} & thallium stress test result coded 4 (meaning not recorded): 31\% to 31\%; thallium stress test result other than thallium stress test result coded 4 (meaning not recorded): 16\% to 50\% & log-odds difference -1.23, OR 0.29 & 2, 240 (small n) & train; test unclear (AUC gain -0.001 [-0.006, 0.000], n=61) & \textcolor{OliveGreen}{yes} & the model has little data on this; on the training data only \\
C39 & interaction: effect of ST depression on exercise (0.0 mm to 2.8 mm) by resting ECG: resting ECG other than a normal resting ECG minus a normal resting ECG; stronger in a normal resting ECG \newline\textit{shape: hinge at 0.4 mm; 57\% of patients same; fallback} & resting ECG other than a normal resting ECG: 25\% to 48\%; a normal resting ECG: 16\% to 50\% & log-odds difference -0.68, OR 0.50 & 121, 121 & train; test agrees (AUC gain -0.002 [-0.016, 0.012], n=61); group difference on test -0.65 [-0.72, -0.58] & \textcolor{OliveGreen}{yes} & clearly \\
C40 & interaction: effect of thallium stress test result (thallium stress test result other than thallium stress test result coded 4 (meaning not recorded) to thallium stress test result coded 4 (meaning not recorded)) by resting ECG: resting ECG other than a normal resting ECG minus a normal resting ECG; stronger in resting ECG other than a normal resting ECG \newline\textit{shape: categorical; highest to lowest: thallium stress test result other than thallium stress test result coded 4 (meaning not recorded), thallium stress test result coded 4 (meaning not recorded)} & resting ECG other than a normal resting ECG: 25\% to 44\%; a normal resting ECG: 20\% to 31\% & log-odds difference 0.14, OR 1.16 & 121, 121, 2 (small n) & train; test unclear (AUC gain -0.002 [-0.016, 0.012], n=61); group difference on test 0.13 [-0.12, 0.35] & \textcolor{OliveGreen}{yes} & the model has little data on this; on the training data only \\
C41 & threshold: ST depression on exercise: fold at about 1.3 mm (starts) \newline\textit{shape: hinge at 1.3 mm; 70\% of patients same; fallback; one point with C41 (median fold 1.3 mm, typical patient 0.4 mm, 10th-90th percentile 0.4 mm to 1.3 mm, wide)} & 16\% to 62\% (typical patient with a fold) & fold 1.28 mm (IQR 0.35 to 1.28); past it log-odds 1.70, OR 5.47; range across patients 1.25 to 2.17 & 242, 109 & train; test unclear (AUC gain 0.009 [-0.008, 0.033], n=61) & \textcolor{OliveGreen}{yes} & on the training data only \\
C42 & subgroup: who this part acts on: 109 of 242 patients & 76\% to 89\% (typical patient this part acts on (risk without the part, then with it)) & share 0.45; log-odds 0.74, OR 2.09; range across patients 0.02 to 1.62 & 109, 133 & train; test unclear (AUC gain 0.006 [-0.014, 0.032], n=61) & \textcolor{OliveGreen}{yes} & on the training data only \\
C43 & structure: its formula reproduces it on 242 training patients (largest difference 1.77636e-15) &  & log-odds 2e-15, OR n/a & 242 & train & \textcolor{OliveGreen}{yes} &  \\
C44 & direction: resting blood pressure: 110 mm Hg to 155 mm Hg (through this part) \newline\textit{shape: hinge at 125 mm Hg; 75\% of patients same; fallback; one point with C52 (median fold 125 mm Hg, typical patient 148 mm Hg, 10th-90th percentile 122 mm Hg to 165 mm Hg, wide)} & 23\% to 17\% & log-odds -0.27, OR 0.77; range across patients -0.48 to -0.08 & 242 & train; test disagrees (AUC gain -0.003 [-0.014, 0.004], n=61) & \textcolor{OliveGreen}{yes} & on the training data only \\
C45 & direction: exercise-induced angina: no chest pain on exercise to chest pain brought on by exercise (through this part) \newline\textit{shape: categorical; highest to lowest: no chest pain on exercise, chest pain brought on by exercise} & 20\% to 18\% & log-odds -0.04, OR 0.96; range across patients -0.08 to 0.00 & 160, 82 & train; test disagrees (AUC gain -0.001 [-0.006, 0.000], n=61) & \textcolor{OliveGreen}{yes} & slightly; on the training data only \\
C46 & direction: chest pain type: chest pain type other than non-anginal chest pain to non-anginal chest pain (through this part) \newline\textit{shape: categorical; highest to lowest: chest pain type other than non-anginal chest pain, non-anginal chest pain} & 20\% to 19\% & log-odds -0.02, OR 0.98; range across patients -0.04 to 0.00 & 172, 70 & train; test unclear (AUC gain 0.000 [0.000, 0.000], n=61) & \textcolor{OliveGreen}{yes} & slightly; on the training data only \\
C47 & direction: thallium stress test result: thallium stress test result other than a fixed defect on the thallium scan to a fixed defect on the thallium scan (through this part) \newline\textit{shape: categorical; highest to lowest: thallium stress test result other than a fixed defect on the thallium scan, a fixed defect on the thallium scan} & 20\% to 19\% & log-odds -0.004, OR 1.00; range across patients -0.007 to 0.00 & 227, 15 (small n) & train; test unclear (AUC gain 0.000 [0.000, 0.000], n=61) & \textcolor{OliveGreen}{yes} & slightly; the model has little data on this; on the training data only \\
C48 & interaction: effect of resting blood pressure (110 mm Hg to 155 mm Hg) by chest pain type: non-anginal chest pain minus chest pain type other than non-anginal chest pain; stronger in non-anginal chest pain \newline\textit{shape: accelerating; 84\% of patients same; fallback; risk saturates near 0\%} & non-anginal chest pain: 3\% to 7\%; chest pain type other than non-anginal chest pain: 15\% to 28\% & log-odds difference 0.06, OR 1.06 & 70, 172 & train; test unclear (AUC gain 0.000 [0.000, 0.000], n=61); group difference on test 0.04 [-0.06, 0.13] & \textcolor{BrickRed}{no} & on the training data only \\
\bottomrule
\end{tabularx}}
\par{\tiny\noindent
\begin{tabularx}{\linewidth}{@{}l>{\raggedright\arraybackslash}X>{\raggedright\arraybackslash}p{0.11\linewidth}>{\raggedright\arraybackslash}p{0.16\linewidth}>{\raggedright\arraybackslash}p{0.05\linewidth}>{\raggedright\arraybackslash}p{0.15\linewidth}l>{\raggedright\arraybackslash}p{0.08\linewidth}@{}}
\toprule
Claim & What is checked & Risk, typical patient & Estimate (log-odds; OR); range across patients & n & Test check & Verified & Words \\
\midrule
C49 & interaction: effect of chest pain type (chest pain type other than non-anginal chest pain to non-anginal chest pain) by exercise-induced angina: chest pain brought on by exercise minus no chest pain on exercise; stronger in chest pain brought on by exercise \newline\textit{shape: categorical; highest to lowest: chest pain type other than non-anginal chest pain, non-anginal chest pain} & chest pain brought on by exercise: 31\% to 7\%; no chest pain on exercise: 20\% to 4\% & log-odds difference -0.001, OR 1.00 & 82, 160 & train; test unclear (AUC gain -0.001 [-0.006, 0.000], n=61); group difference on test -0.008 [-0.02, 0.001] & \textcolor{BrickRed}{no} & on the training data only \\
C50 & interaction: effect of resting blood pressure (110 mm Hg to 155 mm Hg) by exercise-induced angina: chest pain brought on by exercise minus no chest pain on exercise; stronger in no chest pain on exercise \newline\textit{shape: accelerating; 81\% of patients same; fallback} & chest pain brought on by exercise: 25\% to 45\%; no chest pain on exercise: 15\% to 28\% & log-odds difference -0.05, OR 0.95 & 82, 160 & train; test unclear (AUC gain -0.001 [-0.006, 0.000], n=61); group difference on test -0.08 [-0.17, 0.007] & \textcolor{BrickRed}{no} & on the training data only \\
C51 & interaction: effect of resting blood pressure (110 mm Hg to 155 mm Hg) by thallium stress test result: a fixed defect on the thallium scan minus thallium stress test result other than a fixed defect on the thallium scan; stronger in thallium stress test result other than a fixed defect on the thallium scan \newline\textit{shape: accelerating; 76\% of patients same; fallback} & a fixed defect on the thallium scan: 8\% to 17\%; thallium stress test result other than a fixed defect on the thallium scan: 15\% to 28\% & log-odds difference -0.05, OR 0.95 & 15, 227 (small n) & train; test unclear (AUC gain 0.000 [0.000, 0.000], n=61); group difference on test -0.08 [-0.22, 0.12] & \textcolor{BrickRed}{no} & slightly; the model has little data on this; on the training data only \\
C52 & threshold: resting blood pressure: fold at about 125 mm Hg (stops) for patients without chest pain on exercise \newline\textit{shape: hinge at 125 mm Hg; 92\% of patients same; verified; risk saturates near 0\%; one point with C52 (median fold 125 mm Hg, typical patient 122 mm Hg, 10th-90th percentile 122 mm Hg to 165 mm Hg, wide)} & 4\% to 3\% (typical patient with a fold (patients without chest pain on exercise)) & fold 125 mm Hg (IQR 122 to 132); past it log-odds -0.30, OR 0.74; range across patients -0.63 to -0.17 & 143, 78 & train; test disagrees (AUC gain -0.003 [-0.014, 0.004], n=61) & \textcolor{OliveGreen}{yes} & slightly; on the training data only \\
C53 & subgroup: who this part acts on: 145 of 242 patients & 51\% to 68\% (typical patient this part acts on (risk without the part, then with it)) & share 0.60; log-odds 0.52, OR 1.69; range across patients 0.03 to 1.22 & 145, 97 & train; test unclear (AUC gain 0.004 [-0.013, 0.024], n=61) & \textcolor{OliveGreen}{yes} & on the training data only \\
C54 & structure: its formula reproduces it on 242 training patients (largest difference 1.77636e-15) &  & log-odds 2e-15, OR n/a & 242 & train & \textcolor{OliveGreen}{yes} &  \\
C55 & direction: thallium stress test result: thallium stress test result other than a fixed defect on the thallium scan to a fixed defect on the thallium scan (through this part) \newline\textit{shape: categorical; highest to lowest: a fixed defect on the thallium scan, thallium stress test result other than a fixed defect on the thallium scan} & 20\% to 23\% & log-odds 0.18, OR 1.20; range across patients 0.12 to 0.29 & 227, 15 (small n) & train; test unclear (AUC gain 0.002 [0.000, 0.01], n=61) & \textcolor{OliveGreen}{yes} & slightly; the model has little data on this; on the training data only \\
C56 & subgroup: who this part acts on: 234 of 242 patients & 24\% to 20\% (typical patient this part acts on (risk without the part, then with it)) & share 0.97; log-odds -0.23, OR 0.79; range across patients -0.48 to -0.12 & 234, 8 & train; test unclear (AUC gain 0.000 [-0.011, 0.01], n=61) & \textcolor{OliveGreen}{yes} & slightly; on the training data only \\
C57 & structure: its formula reproduces it on 242 training patients (largest difference 2.22045e-16) &  & log-odds 2e-16, OR n/a & 242 & train & \textcolor{OliveGreen}{yes} &  \\
C58 & direction: exercise-induced angina: no chest pain on exercise to chest pain brought on by exercise (through this part) \newline\textit{shape: categorical; highest to lowest: no chest pain on exercise, chest pain brought on by exercise} & 30\% to 25\% (typical patient this part acts on) & log-odds -0.12, OR 0.89; range across patients -0.24 to 0.00 & 160, 82 & train; test disagrees (AUC gain -0.003 [-0.013, 0.000], n=61) & \textcolor{OliveGreen}{yes} & slightly; on the training data only \\
C59 & direction: resting ECG: resting ECG other than signs of left ventricular hypertrophy on the resting ECG to signs of left ventricular hypertrophy on the resting ECG (through this part) \newline\textit{shape: categorical; highest to lowest: signs of left ventricular hypertrophy on the resting ECG, resting ECG other than signs of left ventricular hypertrophy on the resting ECG} & 20\% to 24\% & log-odds 0.16, OR 1.17; range across patients 0.00 to 0.24 & 123, 119 & train; test unclear (AUC gain 0.001 [-0.01, 0.012], n=61) & \textcolor{OliveGreen}{yes} & slightly; on the training data only \\
C60 & interaction: effect of exercise-induced angina (no chest pain on exercise to chest pain brought on by exercise) by resting ECG: signs of left ventricular hypertrophy on the resting ECG minus resting ECG other than signs of left ventricular hypertrophy on the resting ECG; stronger in resting ECG other than signs of left ventricular hypertrophy on the resting ECG \newline\textit{shape: categorical; highest to lowest: chest pain brought on by exercise, no chest pain on exercise} & signs of left ventricular hypertrophy on the resting ECG: 25\% to 35\%; resting ECG other than signs of left ventricular hypertrophy on the resting ECG: 20\% to 31\% & log-odds difference -0.22, OR 0.80 & 119, 123 & train; test agrees (AUC gain 0.001 [-0.01, 0.012], n=61); group difference on test -0.22 [-0.24, -0.20] & \textcolor{OliveGreen}{yes} & slightly \\
\bottomrule
\end{tabularx}}
\par{\tiny\noindent
\begin{tabularx}{\linewidth}{@{}l>{\raggedright\arraybackslash}X>{\raggedright\arraybackslash}p{0.11\linewidth}>{\raggedright\arraybackslash}p{0.16\linewidth}>{\raggedright\arraybackslash}p{0.05\linewidth}>{\raggedright\arraybackslash}p{0.15\linewidth}l>{\raggedright\arraybackslash}p{0.08\linewidth}@{}}
\toprule
Claim & What is checked & Risk, typical patient & Estimate (log-odds; OR); range across patients & n & Test check & Verified & Words \\
\midrule
C61 & subgroup: who this part acts on: 74 of 242 patients & 25\% to 30\% (typical patient this part acts on (risk without the part, then with it)) & share 0.31; log-odds 0.24, OR 1.27; range across patients 0.24 to 0.24 & 74, 168 & train; test disagrees (AUC gain -0.002 [-0.016, 0.006], n=61) & \textcolor{OliveGreen}{yes} & slightly; on the training data only \\
C62 & null: this part barely changes predictions (switched off) &  & max change 0.059 (max $|$change in predicted risk$|$) & 242 & train & \textcolor{BrickRed}{no} &  \\
C63 & structure: its formula reproduces it on 242 training patients (largest difference 0.0) &  & log-odds 0.00, OR n/a & 242 & train & \textcolor{OliveGreen}{yes} &  \\
C64 & direction: fasting blood sugar: fasting blood sugar of 120 mg/dL or less to fasting blood sugar above 120 mg/dL (through this part) \newline\textit{shape: categorical; highest to lowest: fasting blood sugar above 120 mg/dL, fasting blood sugar of 120 mg/dL or less} & 20\% to 21\% & log-odds 0.07, OR 1.08; range across patients 0.07 to 0.07 & 209, 33 & train; test unclear (AUC gain 0.000 [-0.005, 0.004], n=61) & \textcolor{OliveGreen}{yes} & slightly; on the training data only \\
C65 & subgroup: who this part acts on: 33 of 242 patients & 24\% to 25\% (typical patient this part acts on (risk without the part, then with it)) & share 0.14; log-odds 0.07, OR 1.08; range across patients 0.07 to 0.07 & 33, 209 & train; test unclear (AUC gain 0.000 [-0.005, 0.004], n=61) & \textcolor{OliveGreen}{yes} & slightly; on the training data only \\
C66 & null: this part barely changes predictions (switched off) &  & max change 0.019 (max $|$change in predicted risk$|$) & 242 & train & \textcolor{OliveGreen}{yes} &  \\
C67 & structure: its formula reproduces it on 242 training patients (largest difference 0.0) &  & log-odds 0.00, OR n/a & 242 & train & \textcolor{OliveGreen}{yes} &  \\
\bottomrule
\end{tabularx}}
\paragraph{Not supported by the model.}
\begin{itemize}\setlength\itemsep{0pt}\small
\item C16: direction claim about thallium stress test result in the whole model: not supported by the model (15\% of patients have the opposite sign (allowed: 10\%)).
\item C25: direction claim about resting ECG in the whole model: not supported by the model (34\% of patients have the opposite sign (allowed: 10\%)).
\item C28: direction claim about resting ECG in the whole model: not supported by the model (41\% of patients have the opposite sign (allowed: 10\%)).
\item C48: interaction claim about resting blood pressure and chest pain type in part ann\_\allowbreak{}34: not supported by the model (the group difference is below 0.1 log-odds).
\item C49: interaction claim about chest pain type and exercise-induced angina in part ann\_\allowbreak{}34: not supported by the model (the group difference is below 0.1 log-odds; the part's own counterfactual difference has the opposite sign).
\item C50: interaction claim about resting blood pressure and exercise-induced angina in part ann\_\allowbreak{}34: not supported by the model (the group difference is below 0.1 log-odds; the part's own counterfactual difference has the opposite sign).
\item C51: interaction claim about resting blood pressure and thallium stress test result in part ann\_\allowbreak{}34: not supported by the model (the group difference is below 0.1 log-odds; the part's own counterfactual difference has the opposite sign).
\item C62: null claim about resting ECG and exercise-induced angina in part ann\_\allowbreak{}36: not supported by the model (predictions move by up to 0.059).
\end{itemize}

\subsection{Coverage}\label{app:heart-s5-c2:coverage}
Coverage is computed by the current code (the function behind the \texttt{get\_coverage} tool) under these conventions. \emph{Candidate nodes} are all nodes of the explained model state except the output nodes. A candidate node is \emph{covered} when some annotation of the hierarchy (a leaf, or a composition together with its descendants) contains it and all its enabled outgoing edges. \emph{Exact Null} parts, which have no effect on any output for any input, count as covered: nodes with no path to an output over enabled, non-zero-weight edges (disconnected inputs and dead ends), and edges that have zero weight or touch such a node. Compositional coverage is the fraction of required compositions present: every annotation except a single global root containing every output must have a parent; a parentless, childless annotation made only of Null nodes is exempt.

\begin{itemize}
\item \textbf{Null, disconnected input (5):} \texttt{-13} (restecg:1), \texttt{-20} (slope:2), \texttt{-26} (thal:7), \texttt{-3} (sex:1), \texttt{-5} (cp:2).
\item \textbf{Uncovered nodes:} none.
\item \textbf{Uncovered edges:} none.
\item \textbf{Missing compositions:} none.
\end{itemize}

\subsection{Annotation hierarchy}\label{app:heart-s5-c2:hierarchy}
The hierarchy has 1 root; the global root is heart\_\allowbreak{}disease\_\allowbreak{}logit. Tags (A1, A2, \ldots) are used in the figures.
\begingroup\parindent=0pt\parskip=1pt
\hspace*{0.0em}$\triangleright$~A1: heart\_\allowbreak{}disease\_\allowbreak{}logit (composition, \cref{app:heart-s5-c2:ann-heart-disease-logit})\par
\hspace*{1.5em}$\triangleright$~A2: ischaemia\_\allowbreak{}gate (leaf, \cref{app:heart-s5-c2:ann-ischaemia-gate})\par
\hspace*{1.5em}$\triangleright$~A3: bp\_\allowbreak{}offset\_\allowbreak{}gate (leaf, \cref{app:heart-s5-c2:ann-bp-offset-gate})\par
\hspace*{1.5em}$\triangleright$~A4: ischaemia\_\allowbreak{}damping\_\allowbreak{}thal6 (leaf, \cref{app:heart-s5-c2:ann-ischaemia-damping-thal6})\par
\hspace*{1.5em}$\triangleright$~A5: lvh\_\allowbreak{}without\_\allowbreak{}exercise\_\allowbreak{}angina (leaf, \cref{app:heart-s5-c2:ann-lvh-without-exercise-angina})\par
\hspace*{1.5em}$\triangleright$~A6: high\_\allowbreak{}fasting\_\allowbreak{}sugar (leaf, \cref{app:heart-s5-c2:ann-high-fasting-sugar})\par
\endgroup

\subsection{Network figures}\label{app:heart-s5-c2:figures}
\begin{figure}[htbp]
\centering
\input{\explaneatcertdir{heart-s5-c2}heart-s5-c2-fig-full.tex}
\caption{heart-s5-c2: full view of the explained model state.}
\label{app:heart-s5-c2:fig-full}
\end{figure}
\begin{figure}[htbp]
\centering
\input{\explaneatcertdir{heart-s5-c2}heart-s5-c2-fig-level-0.tex}
\caption{heart-s5-c2: collapsed view at hierarchy level 0 (fully recomposed).}
\label{app:heart-s5-c2:fig-level-0}
\end{figure}
\begin{figure}[htbp]
\centering
\input{\explaneatcertdir{heart-s5-c2}heart-s5-c2-fig-level-1.tex}
\caption{heart-s5-c2: collapsed view at hierarchy level 1.}
\label{app:heart-s5-c2:fig-level-1}
\end{figure}

\subsection{Refactoring operations}\label{app:heart-s5-c2:operations}
The explained model state is the phenotype after these operations, replayed in order. All are function-preserving; the largest change in the model output over the 242 training rows is $2.22\times 10^{-16}$ (tolerance $1.00\times 10^{-9}$).
\begin{itemize}
\item \#26: split \texttt{-12} into \texttt{-12\_\allowbreak{}a}, \texttt{-12\_\allowbreak{}b}.
\item \#27: split \texttt{-17} into \texttt{-17\_\allowbreak{}a}, \texttt{-17\_\allowbreak{}b}, \texttt{-17\_\allowbreak{}c}.
\item \#28: split \texttt{-25} into \texttt{-25\_\allowbreak{}a}, \texttt{-25\_\allowbreak{}b}, \texttt{-25\_\allowbreak{}c}.
\item \#29: split \texttt{-8} into \texttt{-8\_\allowbreak{}a}, \texttt{-8\_\allowbreak{}b}.
\item \#30: split \texttt{-6} into \texttt{-6\_\allowbreak{}a}, \texttt{-6\_\allowbreak{}b}.
\item \#31: identity node \texttt{identity\_\allowbreak{}290} before \texttt{290} (1 connection(s)).
\item \#32: identity node \texttt{identity\_\allowbreak{}1416} before \texttt{1416} (1 connection(s)).
\end{itemize}

\subsection{Annotations}\label{app:heart-s5-c2:annotations}
\subsubsection{\texorpdfstring{A1: heart\_\allowbreak{}disease\_\allowbreak{}logit}{A1: heart disease logit}}\label{app:heart-s5-c2:ann-heart-disease-logit}
% annotation id ann_38, operation seq 38
\textbf{Kind:} Composition annotation (depth 0); children A2 (\cref{app:heart-s5-c2:ann-ischaemia-gate}), A3 (\cref{app:heart-s5-c2:ann-bp-offset-gate}), A4 (\cref{app:heart-s5-c2:ann-ischaemia-damping-thal6}), A5 (\cref{app:heart-s5-c2:ann-lvh-without-exercise-angina}), A6 (\cref{app:heart-s5-c2:ann-high-fasting-sugar}).

\begin{itemize}
\item \textbf{Entry nodes:} \texttt{-1} (age), \texttt{-2} (sex:0), \texttt{-3} (sex:1), \texttt{-4} (cp:1), \texttt{-5} (cp:2), \texttt{-6\_\allowbreak{}a} (cp:3\_\allowbreak{}a), \texttt{-7} (cp:4), \texttt{-8\_\allowbreak{}a} (trestbps\_\allowbreak{}a), \texttt{-9} (chol), \texttt{-10} (fbs:0), \texttt{-12\_\allowbreak{}a} (restecg:0\_\allowbreak{}a), \texttt{-13} (restecg:1), \texttt{-15} (thalach), \texttt{-16} (exang:0), \texttt{-17\_\allowbreak{}a} (exang:1\_\allowbreak{}a), \texttt{-19} (slope:1), \texttt{-20} (slope:2), \texttt{-21} (slope:3), \texttt{-22} (ca), \texttt{-23} (thal:3), \texttt{-25\_\allowbreak{}a} (thal:6\_\allowbreak{}a), \texttt{-26} (thal:7).
\item \textbf{Exit nodes:} \texttt{0}.
\item \textbf{Internal nodes:} none.
\item \textbf{Edges (weight):} \texttt{-1}$\to$\texttt{0}~(0.126), \texttt{-2}$\to$\texttt{0}~(-1.06), \texttt{-4}$\to$\texttt{0}~(-0.489), \texttt{-7}$\to$\texttt{0}~(0.357), \texttt{-9}$\to$\texttt{0}~(0.522), \texttt{-10}$\to$\texttt{0}~(0.264), \texttt{-15}$\to$\texttt{0}~(-0.226), \texttt{-16}$\to$\texttt{0}~(0.136), \texttt{-19}$\to$\texttt{0}~(-0.921), \texttt{-21}$\to$\texttt{0}~(-0.37), \texttt{-22}$\to$\texttt{0}~(1.47), \texttt{-23}$\to$\texttt{0}~(-0.423), \texttt{204}$\to$\texttt{0}~(0.135), \texttt{290}$\to$\texttt{0}~(0.252), \texttt{1416}$\to$\texttt{0}~(-0.118), \texttt{1611}$\to$\texttt{0}~(0.729), \texttt{3884}$\to$\texttt{0}~(0.0482), \texttt{-12\_\allowbreak{}a}$\to$\texttt{0}~(-0.158), \texttt{-17\_\allowbreak{}a}$\to$\texttt{0}~(0.471), \texttt{-25\_\allowbreak{}a}$\to$\texttt{0}~(-0.413), \texttt{-8\_\allowbreak{}a}$\to$\texttt{0}~(0.498), \texttt{-6\_\allowbreak{}a}$\to$\texttt{0}~(-0.444), \texttt{204}$\to$\texttt{identity\_\allowbreak{}290}~(1.14), \texttt{204}$\to$\texttt{identity\_\allowbreak{}1416}~(0.841).
\item \textbf{With its descendants:} 41 nodes.
\item \textbf{Covers} (coverage, own and as part of the hierarchy): 40 candidate nodes.
\end{itemize}

\textbf{Hypothesis.} Root: p = sigmoid(z), z = -0.487 + LINEAR CORE + 0.136*ischaemia\_\allowbreak{}gate + 0.252*bp\_\allowbreak{}offset\_\allowbreak{}gate - 0.118*\allowbreak{}ischaemia\_\allowbreak{}damping\_\allowbreak{}thal6 + 0.729*\allowbreak{}lvh\_\allowbreak{}without\_\allowbreak{}exercise\_\allowbreak{}angina + 0.048*\allowbreak{}high\_\allowbreak{}fasting\_\allowbreak{}sugar. LINEAR CORE (17 direct input-$>$output edges, z-scored inputs): ca +1.473, sex:0 -1.062, slope:1 -0.921, chol +0.522, trestbps +0.498, cp:1 -0.489, exang:1 +0.471, cp:3 -0.444, thal:3 -0.423, thal:6 -0.413, slope:3 -0.370, cp:4 +0.357, fbs:0 +0.264, thalach -0.226, restecg:0 -0.159, exang:0 +0.136, age +0.126. Five inputs are unconnected (sex:1, cp:2, restecg:1, slope:2, thal:7; zero-effect entries of this root). Oldpeak, thal:4, restecg:2 and fbs:1 have no direct edge and act only through the children. COMPOSITION: additive readouts, with a serial fan-out inside. ischaemia\_\allowbreak{}gate (oldpeak above a resting-ECG-dependent threshold) is read out directly and also feeds bp\_\allowbreak{}offset\_\allowbreak{}gate (the ischaemia score minus a resting-BP offset, the largest nonlinear readout) and ischaemia\_\allowbreak{}damping\_\allowbreak{}thal6 (a near-linear term that takes back part of the ischaemia score and adds a fixed-defect term); identity nodes identity\_\allowbreak{}290 and identity\_\allowbreak{}1416 carry those two links so each gate is its own leaf. Two independent indicators add to z: lvh\_\allowbreak{}without\_\allowbreak{}exercise\_\allowbreak{}angina (+0.236 for LVH on the resting ECG without exercise angina) and high\_\allowbreak{}fasting\_\allowbreak{}sugar (+0.074, negligible). EVIDENCE (train, n=242): full model AUC 0.9389, acc 0.880 (compute\_\allowbreak{}performance)\allowbreak{}. Linear core alone AUC 0.933; logit SD 3.19 (core) vs 0.57 (all hidden readouts). Removing the ischaemia gate with its downstream effects gives 0.932; readout-only ablations: ischaemia 0.938, BP 0.933, damping 0.938, LVH 0.939, fbs 0.939. Treating all five ReLUs as identities gives 0.917 (the gating matters); one least-squares linear fit of all hidden readouts on their inputs gives 0.938 (R2 0.84); an L2 logistic regression on all 26 inputs gives 0.936. Mean $|$SHAP$|$ (train, all 242): ca 0.145, sex:0 0.109, slope:1 0.108, thal:3 0.050, oldpeak 0.047, cp:3 0.045, exang:1 0.044, chol 0.043. Interpretation: close to logistic regression plus one ischaemia correction: ST depression raises risk only above a threshold that is lower with a normal resting ECG, more steeply when resting BP is low; LVH without exercise angina adds a little. Counter-directional signs reported, not explained: restecg:0 positive in the ischaemia gate, trestbps negative in the BP gate, thal:6 (fixed defect) negative directly, fbs:0 +0.26 directly. Same core design as the 10-seed heart-s5 record (key 9330) and the sibling cluster lead heart-s5-c1 (ischaemia gate -$>$ resting-BP gate), with a damping node, an LVH indicator and a fasting-sugar indicator instead of c1's age gate.

\textbf{Closed form} (children as named functions):
{\small\begingroup\allowdisplaybreaks
\begin{align*}
y_{0} &= \sigma\bigl(z_{0}\bigr) \\
z_{0} &= -0.487 + 0.136 x_{exang:0} + 0.522 x_{chol} + 0.471 x_{exang:1\_a} + 0.126 x_{age} \\
&\qquad + 0.357 x_{cp:4} \\
&\qquad + 0.252 \operatorname{bp\_offset\_gate}\bigl(1.14 \operatorname{ischaemia\_gate}\bigl(x_{restecg:0\_b}, \\
&\qquad x_{oldpeak}, x_{thal:4}\bigr), x_{trestbps\_b}, x_{cp:3\_b}, x_{exang:1\_b}, \\
&\qquad x_{thal:6\_b}\bigr) + 0.264 x_{fbs:0} \\
&\qquad + 0.135 \operatorname{ischaemia\_gate}\bigl(x_{restecg:0\_b}, x_{oldpeak}, x_{thal:4}\bigr) \\
&\qquad + 0.729 \operatorname{lvh\_without\_exercise\_angina}\bigl(x_{restecg:2}, x_{exang:1\_c}\bigr) \\
&\qquad + 1.47 x_{ca} + 0.0482 \operatorname{high\_fasting\_sugar}\bigl(x_{fbs:1}\bigr) \\
&\qquad + 0.497 x_{trestbps\_a} - 1.06 x_{sex:0} - 0.423 x_{thal:3} - 0.444 x_{cp:3\_a} \\
&\qquad - 0.158 x_{restecg:0\_a} - 0.37 x_{slope:3} - 0.413 x_{thal:6\_a} - 0.226 x_{thalach} \\
&\qquad - 0.921 x_{slope:1} \\
&\qquad - 0.118 \operatorname{ischaemia\_damping\_thal6}\bigl(0.841 \operatorname{ischaemia\_gate}\bigl(x_{restecg:0\_b}, \\
&\qquad x_{oldpeak}, x_{thal:4}\bigr), x_{thal:6\_c}\bigr) - 0.489 x_{cp:1}
\end{align*}\endgroup}

Expanded over the children:
{\small\begingroup\allowdisplaybreaks
\begin{align*}
y_{0} &= \sigma\bigl(z_{0}\bigr) \\
z_{0} &= -0.487 + 0.136 x_{exang:0} + 0.522 x_{chol} + 0.471 x_{exang:1\_a} + 0.126 x_{age} \\
&\qquad + 0.357 x_{cp:4} + 0.252 \operatorname{relu}\bigl(-0.447 + 1.14 \operatorname{relu}\bigl(0.385 \\
&\qquad + 0.992 x_{restecg:0\_b} + 2.54 x_{oldpeak} - 0.876 x_{thal:4}\bigr) - 0.0688 x_{cp:3\_b} \\
&\qquad - 0.153 x_{exang:1\_b} - 0.00683 x_{thal:6\_b} - 0.773 x_{trestbps\_b}\bigr) + 0.264 x_{fbs:0} \\
&\qquad + 0.135 \operatorname{relu}\bigl(0.385 + 0.992 x_{restecg:0\_b} + 2.54 x_{oldpeak} \\
&\qquad - 0.876 x_{thal:4}\bigr) + 0.729 \operatorname{relu}\bigl(-1.26 + 0.887 x_{restecg:2} \\
&\qquad - 0.953 x_{exang:1\_c}\bigr) + 1.47 x_{ca} + 0.0482 \operatorname{relu}\bigl(-0.98 \\
&\qquad + 1.0 x_{fbs:1}\bigr) + 0.497 x_{trestbps\_a} - 1.06 x_{sex:0} - 0.423 x_{thal:3} \\
&\qquad - 0.444 x_{cp:3\_a} - 0.158 x_{restecg:0\_a} - 0.37 x_{slope:3} - 0.413 x_{thal:6\_a} \\
&\qquad - 0.226 x_{thalach} - 0.921 x_{slope:1} - 0.118 \operatorname{relu}\bigl(0.853 \\
&\qquad + 0.841 \operatorname{relu}\bigl(0.385 + 0.992 x_{restecg:0\_b} + 2.54 x_{oldpeak} \\
&\qquad - 0.876 x_{thal:4}\bigr) - 0.594 x_{thal:6\_c}\bigr) - 0.489 x_{cp:1}
\end{align*}\endgroup}

where $y_{0}$ is node \texttt{0}; $z_{i}$, where shown, is the pre-activation of a sigmoid output $y_{i}$; weights and biases are shown to 3 significant figures.
 On the training split (n = 242), the displayed closed form differs from the forward pass by at most 0.00243, and the exact function by at most $2.22\times 10^{-16}$.

\textbf{Evidence} (2 items; full payloads in the JSON record):
\begin{itemize}
\item \textsc{analytical}, \texttt{analytical\_\allowbreak{}formula}; data split: n/a (analytical). Root closed form (get\_\allowbreak{}formula force=true); output is sigmoid.
\item \textsc{empirical}, \texttt{empirical\_\allowbreak{}viz}; data split: train. Train-split performance, ablations, linear replacements and SHAP.
\end{itemize}

\textbf{Authorship:} drafted by AI-assisted draft on 2026-11-10; verified by nobody (unverified).

\subsubsection{\texorpdfstring{A2: ischaemia\_\allowbreak{}gate}{A2: ischaemia gate}}\label{app:heart-s5-c2:ann-ischaemia-gate}
% annotation id ann_33, operation seq 33
\textbf{Kind:} Leaf annotation (depth 1); parent A1 (\cref{app:heart-s5-c2:ann-heart-disease-logit}).

\begin{itemize}
\item \textbf{Entry nodes:} \texttt{-12\_\allowbreak{}b} (restecg:0\_\allowbreak{}b), \texttt{-18} (oldpeak), \texttt{-24} (thal:4).
\item \textbf{Exit nodes:} \texttt{204}.
\item \textbf{Internal nodes:} none.
\item \textbf{Edges (weight):} \texttt{-12\_\allowbreak{}b}$\to$\texttt{204}~(0.992), \texttt{-18}$\to$\texttt{204}~(2.54), \texttt{-24}$\to$\texttt{204}~(-0.876).
\item \textbf{Covers} (coverage, own and as part of the hierarchy): 3 candidate nodes.
\end{itemize}

\textbf{Hypothesis.} Ischaemia-evidence gate (raises risk). 204 = ReLU(0.385 + 2.536*oldpeak + 0.992*restecg:0 - 0.876*thal:4). Hypothesis: a one-sided ST-depression score that opens above an oldpeak threshold set by the resting ECG (z-oldpeak -0.57 with a normal ECG, +0.21 without) and then rises with oldpeak; it adds 0.136 per unit to the logit directly, a further 0.286 per unit through the BP gate while that gate is open, and -0.099 per unit through the damping node, so about 0.32 logit per unit in all. EVIDENCE (train, n=242): active for 45\% (51\% with a normal ECG, 40\% without); corr with oldpeak 0.92; disease rate 60\% when active vs 32\% when inactive. Mean $|$sensitivity$|$ oldpeak 1.14 $>$ restecg:0 0.45 $>$ thal:4 0.39 (thal:4 has 2 patients). This is the main nonlinear effect: removing it with its downstream effects lowers train AUC 0.939 -$>$ 0.932 (below the linear-core level 0.933); an affine replacement recovers 0.938. Supported. Oldpeak and thal:4 have no direct edge and act only here. The positive restecg:0 weight (a normal ECG raises the score; restecg:0's direct edge is -0.16) is counter-intuitive and is reported, not explained.

\textbf{Closed form:}
{\small\begingroup\allowdisplaybreaks
\begin{align*}
y_{0} &= \operatorname{relu}\bigl(0.385 + 0.992 x_{restecg:0\_b} + 2.54 x_{oldpeak} \\
&\qquad - 0.876 x_{thal:4}\bigr)
\end{align*}\endgroup}

where $y_{0}$ is node \texttt{204}; $z_{i}$, where shown, is the pre-activation of a sigmoid output $y_{i}$; weights and biases are shown to 3 significant figures.
 On the training split (n = 242), the displayed closed form differs from the forward pass by at most 0.0173, and the exact function by at most 0.

\textbf{Evidence} (2 items; full payloads in the JSON record):
\begin{itemize}
\item \textsc{analytical}, \texttt{analytical\_\allowbreak{}formula}; data split: n/a (analytical); thresholds: thresholds = z-oldpeak at which 204 opens: -0.57 with a normal resting E.... Closed form from get\_\allowbreak{}formula (force) with readouts and thresholds.
\item \textsc{empirical}, \texttt{empirical\_\allowbreak{}viz} (sensitivity (max\_\allowbreak{}samples 100000)); data split: train. Train-split sensitivity, activity, outcome rates and ablation.
\end{itemize}

\textbf{Authorship:} drafted by AI-assisted draft on 2026-11-10; verified by nobody (unverified).

\subsubsection{\texorpdfstring{A3: bp\_\allowbreak{}offset\_\allowbreak{}gate}{A3: bp offset gate}}\label{app:heart-s5-c2:ann-bp-offset-gate}
% annotation id ann_34, operation seq 34
\textbf{Kind:} Leaf annotation (depth 1); parent A1 (\cref{app:heart-s5-c2:ann-heart-disease-logit}).

\begin{itemize}
\item \textbf{Entry nodes:} \texttt{identity\_\allowbreak{}290}, \texttt{-8\_\allowbreak{}b} (trestbps\_\allowbreak{}b), \texttt{-6\_\allowbreak{}b} (cp:3\_\allowbreak{}b), \texttt{-17\_\allowbreak{}b} (exang:1\_\allowbreak{}b), \texttt{-25\_\allowbreak{}b} (thal:6\_\allowbreak{}b).
\item \textbf{Exit nodes:} \texttt{290}.
\item \textbf{Internal nodes:} none.
\item \textbf{Edges (weight):} \texttt{identity\_\allowbreak{}290}$\to$\texttt{290}~(1), \texttt{-8\_\allowbreak{}b}$\to$\texttt{290}~(-0.773), \texttt{-6\_\allowbreak{}b}$\to$\texttt{290}~(-0.0688), \texttt{-17\_\allowbreak{}b}$\to$\texttt{290}~(-0.153), \texttt{-25\_\allowbreak{}b}$\to$\texttt{290}~(-0.00683).
\item \textbf{Covers} (coverage, own and as part of the hierarchy): 5 candidate nodes.
\end{itemize}

\textbf{Hypothesis.} Blood-pressure offset gate on the ischaemia score (raises risk). 290 = ReLU(-0.447 + 1.137*ischaemia\_\allowbreak{}gate - 0.773*trestbps - 0.153*exang:1 - 0.069*cp:3 - 0.007*thal:6), read out +0.252. Hypothesis: passes the ischaemia score on, minus an offset that grows with resting blood pressure (and slightly with exercise angina and non-anginal pain). With the ischaemia gate closed it still opens for low resting blood pressure, so within this gate lower resting BP raises risk (0.252*0.773 = 0.195 per SD while open), partly offsetting trestbps's direct +0.50 edge. EVIDENCE (train, n=242): active for 60\% (87\% with the ischaemia gate open, 38\% without); with the ischaemia gate open, active for 100\%, 85\% and 77\% of the low, middle and high resting-BP groups. Disease rate 50\% when active vs 35\% when inactive. Mean $|$sensitivity$|$ ischaemia score 0.60 $>$ trestbps 0.46 $>$ exang:1 0.09 $>$ cp:3 0.04 $>$ thal:6 0.004. It is the largest nonlinear readout (logit SD 0.50): readout ablation lowers train AUC 0.939 -$>$ 0.933; an affine replacement fits it well (R2 0.97, AUC 0.939), so its nonlinearity matters little. Supported. The negative BP term inside a risk-raising gate is reported, not explained.

\textbf{Closed form:}
{\small\begingroup\allowdisplaybreaks
\begin{align*}
y_{0} &= \operatorname{relu}\bigl(-0.447 + 1.0 x_{identity\_290} - 0.0688 x_{cp:3\_b} \\
&\qquad - 0.153 x_{exang:1\_b} - 0.00683 x_{thal:6\_b} - 0.773 x_{trestbps\_b}\bigr)
\end{align*}\endgroup}

where $y_{0}$ is node \texttt{290}; $z_{i}$, where shown, is the pre-activation of a sigmoid output $y_{i}$; weights and biases are shown to 3 significant figures.
 On the training split (n = 242), the displayed closed form differs from the forward pass by at most 0.00119, and the exact function by at most $1.78\times 10^{-15}$.

\textbf{Evidence} (2 items; full payloads in the JSON record):
\begin{itemize}
\item \textsc{analytical}, \texttt{analytical\_\allowbreak{}formula}; data split: n/a (analytical). Closed form from get\_\allowbreak{}formula (force) with readouts and thresholds.
\item \textsc{empirical}, \texttt{empirical\_\allowbreak{}viz} (sensitivity (max\_\allowbreak{}samples 100000)); data split: train. Train-split sensitivity, activity, outcome rates and ablation.
\end{itemize}

\textbf{Authorship:} drafted by AI-assisted draft on 2026-11-10; verified by nobody (unverified).

\subsubsection{\texorpdfstring{A4: ischaemia\_\allowbreak{}damping\_\allowbreak{}thal6}{A4: ischaemia damping thal6}}\label{app:heart-s5-c2:ann-ischaemia-damping-thal6}
% annotation id ann_35, operation seq 35
\textbf{Kind:} Leaf annotation (depth 1); parent A1 (\cref{app:heart-s5-c2:ann-heart-disease-logit}).

\begin{itemize}
\item \textbf{Entry nodes:} \texttt{identity\_\allowbreak{}1416}, \texttt{-25\_\allowbreak{}c} (thal:6\_\allowbreak{}c).
\item \textbf{Exit nodes:} \texttt{1416}.
\item \textbf{Internal nodes:} none.
\item \textbf{Edges (weight):} \texttt{identity\_\allowbreak{}1416}$\to$\texttt{1416}~(1), \texttt{-25\_\allowbreak{}c}$\to$\texttt{1416}~(-0.594).
\item \textbf{Covers} (coverage, own and as part of the hierarchy): 2 candidate nodes.
\end{itemize}

\textbf{Hypothesis.} Ischaemia damping with a fixed-defect term (near-linear, small). 1416 = ReLU(0.853 + 0.841*ischaemia\_\allowbreak{}gate - 0.594*thal:6), read out -0.118. Hypothesis: almost always open, so it acts as a linear correction: it subtracts 0.099 per unit of the ischaemia score from the logit (damping the ischaemia gate's direct +0.136) and adds 0.070 per SD of thal:6 (about +0.29 logit for a fixed thallium defect, against thal:6's direct -0.41 edge). It closes only for fixed-defect patients whose ischaemia score is below 1.73. EVIDENCE (train, n=242): active for 97\% (closed for 8 of the 15 fixed-defect patients). Mean $|$sensitivity$|$ ischaemia score 0.97, thal:6 0.57. Readout ablation: train AUC 0.939 -$>$ 0.938; affine replacement R2 0.99, AUC 0.939. Supported; effectively linear.

\textbf{Closed form:}
{\small\begingroup\allowdisplaybreaks
\begin{align*}
y_{0} &= \operatorname{relu}\bigl(0.853 + 1.0 x_{identity\_1416} - 0.594 x_{thal:6\_c}\bigr)
\end{align*}\endgroup}

where $y_{0}$ is node \texttt{1416}; $z_{i}$, where shown, is the pre-activation of a sigmoid output $y_{i}$; weights and biases are shown to 3 significant figures.
 On the training split (n = 242), the displayed closed form differs from the forward pass by at most 0.00134, and the exact function by at most $2.22\times 10^{-16}$.

\textbf{Evidence} (2 items; full payloads in the JSON record):
\begin{itemize}
\item \textsc{analytical}, \texttt{analytical\_\allowbreak{}formula}; data split: n/a (analytical). Closed form from get\_\allowbreak{}formula (force) with readouts and thresholds.
\item \textsc{empirical}, \texttt{empirical\_\allowbreak{}viz} (sensitivity (max\_\allowbreak{}samples 100000)); data split: train. Train-split sensitivity, activity, outcome rates and ablation.
\end{itemize}

\textbf{Authorship:} drafted by AI-assisted draft on 2026-11-10; verified by nobody (unverified).

\subsubsection{\texorpdfstring{A5: lvh\_\allowbreak{}without\_\allowbreak{}exercise\_\allowbreak{}angina}{A5: lvh without exercise angina}}\label{app:heart-s5-c2:ann-lvh-without-exercise-angina}
% annotation id ann_36, operation seq 36
\textbf{Kind:} Leaf annotation (depth 1); parent A1 (\cref{app:heart-s5-c2:ann-heart-disease-logit}).

\begin{itemize}
\item \textbf{Entry nodes:} \texttt{-14} (restecg:2), \texttt{-17\_\allowbreak{}c} (exang:1\_\allowbreak{}c).
\item \textbf{Exit nodes:} \texttt{1611}.
\item \textbf{Internal nodes:} none.
\item \textbf{Edges (weight):} \texttt{-14}$\to$\texttt{1611}~(0.887), \texttt{-17\_\allowbreak{}c}$\to$\texttt{1611}~(-0.953).
\item \textbf{Covers} (coverage, own and as part of the hierarchy): 2 candidate nodes.
\end{itemize}

\textbf{Hypothesis.} LVH without exercise angina (AND indicator, raises risk, small). 1611 = ReLU(-1.261 + 0.887*restecg:2 - 0.953*exang:1), read out +0.729. Hypothesis: a binary AND: open (constant 0.324, +0.236 logit) only for patients with left ventricular hypertrophy on the resting ECG and no exercise-induced angina, zero otherwise. EVIDENCE (train, n=242): active for exactly the 74 patients with LVH and no exercise angina; their disease rate is 41\%, against 20\% for patients with neither. Mean $|$sensitivity$|$ exang:1 0.29, restecg:2 0.27. Readout ablation: train AUC 0.939 -$>$ 0.939 (-0.0003). restecg:2 acts only here. Supported; near-negligible for ranking.

\textbf{Closed form:}
{\small\begingroup\allowdisplaybreaks
\begin{align*}
y_{0} &= \operatorname{relu}\bigl(-1.26 + 0.887 x_{restecg:2} - 0.953 x_{exang:1\_c}\bigr)
\end{align*}\endgroup}

where $y_{0}$ is node \texttt{1611}; $z_{i}$, where shown, is the pre-activation of a sigmoid output $y_{i}$; weights and biases are shown to 3 significant figures.
 On the training split (n = 242), the displayed closed form differs from the forward pass by at most $2.52\times 10^{-4}$, and the exact function by at most 0.

\textbf{Evidence} (2 items; full payloads in the JSON record):
\begin{itemize}
\item \textsc{analytical}, \texttt{analytical\_\allowbreak{}formula}; data split: n/a (analytical). Closed form from get\_\allowbreak{}formula (force) with readouts and thresholds.
\item \textsc{empirical}, \texttt{empirical\_\allowbreak{}viz} (sensitivity (max\_\allowbreak{}samples 100000)); data split: train. Train-split sensitivity, activity, outcome rates and ablation.
\end{itemize}

\textbf{Authorship:} drafted by AI-assisted draft on 2026-11-10; verified by nobody (unverified).

\subsubsection{\texorpdfstring{A6: high\_\allowbreak{}fasting\_\allowbreak{}sugar}{A6: high fasting sugar}}\label{app:heart-s5-c2:ann-high-fasting-sugar}
% annotation id ann_37, operation seq 37
\textbf{Kind:} Leaf annotation (depth 1); parent A1 (\cref{app:heart-s5-c2:ann-heart-disease-logit}).

\begin{itemize}
\item \textbf{Entry nodes:} \texttt{-11} (fbs:1).
\item \textbf{Exit nodes:} \texttt{3884}.
\item \textbf{Internal nodes:} none.
\item \textbf{Edges (weight):} \texttt{-11}$\to$\texttt{3884}~(1).
\item \textbf{Covers} (coverage, own and as part of the hierarchy): 1 candidate node.
\end{itemize}

\textbf{Hypothesis.} High fasting blood sugar indicator (negligible). 3884 = ReLU(-0.980 + fbs:1), read out +0.048: +0.074 logit for the 33 patients with fasting blood sugar $>$ 120 mg/dl. Hypothesis: a small additive term for high fasting sugar (exactly linear on the two values of fbs:1); together with fbs:0's direct +0.26 edge, high fasting sugar overall lowers the logit (fbs:0 is the normal-sugar flag). EVIDENCE (train, n=242): active for exactly the 33 high-sugar patients; disease rate 45\% vs 44\%. Readout ablation changes train AUC by less than 0.0001; affine replacement R2 1.0. Supported; negligible.

\textbf{Closed form:}
{\small\begingroup\allowdisplaybreaks
\begin{align*}
y_{0} &= \operatorname{relu}\bigl(-0.98 + 1.0 x_{fbs:1}\bigr)
\end{align*}\endgroup}

where $y_{0}$ is node \texttt{3884}; $z_{i}$, where shown, is the pre-activation of a sigmoid output $y_{i}$; weights and biases are shown to 3 significant figures.
 On the training split (n = 242), the displayed closed form differs from the forward pass by at most $3.12\times 10^{-4}$, and the exact function by at most 0.

\textbf{Evidence} (2 items; full payloads in the JSON record):
\begin{itemize}
\item \textsc{analytical}, \texttt{analytical\_\allowbreak{}formula}; data split: n/a (analytical). Closed form from get\_\allowbreak{}formula (force) with readouts and thresholds.
\item \textsc{empirical}, \texttt{empirical\_\allowbreak{}viz} (sensitivity (max\_\allowbreak{}samples 100000)); data split: train. Train-split sensitivity, activity, outcome rates and ablation.
\end{itemize}

\textbf{Authorship:} drafted by AI-assisted draft on 2026-11-10; verified by nobody (unverified).

\subsection{Audit against the definitions}\label{app:heart-s5-c2:audit}
Automatic findings on the stored explanation (codes as in the JSON record's \texttt{issues}); they flag what an examiner should check, not verdicts.
\begin{itemize}
\item \textsc{warning} \texttt{unverified}: No human verification is recorded for this explanation.
\end{itemize}
\end{drafted}
